\documentclass[11pt]{article}
\usepackage[margin=1in]{geometry}
\usepackage{amsmath,amssymb,amsthm}
\usepackage{hyperref}
\usepackage{listings}
\usepackage[numbers,sort&compress]{natbib}

\newtheorem{theorem}{Theorem}
\newtheorem{proposition}{Proposition}
\newtheorem{definition}{Definition}

\lstset{basicstyle=\ttfamily\small,columns=fullflexible,breaklines=true,
  frame=single,framesep=4pt,xleftmargin=4pt}

\title{A Machine-Checked Core for Set Theoretic Estimation}
\author{Project \texttt{thejeb}}
\date{\today}

\begin{document}
\maketitle

\begin{abstract}
We report a first, machine-checked formalization in Lean~4 / Mathlib of
the core of set theoretic estimation (STE) as founded by
\citet{combettes1993} and recast over topological spaces for decryption
by \citet{carlson2012}. The formalization covers the feasibility set and
its lattice-theoretic properties, the consistent / fair / ideal
taxonomy, information monotonicity, an inconsistency-based
invalid-information detector, closedness and compactness-based existence
of estimates on topological spaces, convexity of the feasibility set,
a skeleton of Carlson's STE-decryption reduction, and --- new in this
revision --- Carlson's key-counting Lemma~4.1, the cipher-reduction
chain (Theorem~5.2, Lemma~5.1 / Corollary~5.2, Corollary~5.3) and a
deterministic form of asymptotic unicity. Every theorem
stated here is verified by continuous integration on each push; nothing
is asserted proven otherwise.
\end{abstract}

\section{Introduction}
Set theoretic estimation abandons the search for an optimum in favor of
\emph{feasibility}: an estimate is acceptable exactly when it is
consistent with every available piece of information
\citep{combettes1993}. This shifts the mathematical content from analysis
of an objective function to the set algebra and topology of a family of
constraint sets and their intersection. That content is a natural target
for a proof assistant. We mechanize it in Lean~4 on Mathlib (pinned
\texttt{v4.32.0}), adopting a strict discipline: a statement counts as
established only when it type-checks in our library under the project's
CI. This paper is the human-readable companion to that library; every
displayed result names its Lean identifier.

\section{The formal model}
Let $\Xi$ be the solution space and $I$ an index set of information
pieces. A \emph{property set family} is $S:I\to\mathcal P(\Xi)$; $S_i$
holds the estimates consistent with information $i$. We mechanize the
crisp case (Combettes's membership function $\Psi_i$ being
$\{0,1\}$-valued, so $S_i$ is a subset).

\begin{definition}[feasibility set; \texttt{STE.feasibilitySet}]
$S \;=\; \bigcap_{i\in I} S_i$ (Combettes Eq.~4). A point of $S$ is a
\emph{set theoretic estimate}.
\end{definition}

\begin{lstlisting}
def feasibilitySet (S : I -> Set Xi) : Set Xi := iInter i, S i
theorem mem_feasibilitySet : a in feasibilitySet S <-> forall i, a in S i
\end{lstlisting}

\section{Results (all machine-checked)}

\paragraph{Refinement and monotonicity.}
The feasibility set refines each constraint
(\texttt{feasibilitySet\_subset}: $S\subseteq S_j$). Writing
$S_J=\bigcap_{i\in J}S_i$ for the \emph{partial} feasibility set,
enforcing more information shrinks it:
\begin{proposition}[\texttt{partialFeasibilitySet\_antitone}]
$J\subseteq K \implies S_K \subseteq S_J$.
\end{proposition}
\noindent This is the STE reading of ``more information never enlarges
the acceptable set,'' and it descends to Carlson's key search as
\texttt{feasibleKeys\_antitone}.

\paragraph{Consistency, fairness, invalid information.}
Call the family \emph{fair} for $h$ when $h\in S_i$ for all $i$
(\texttt{STE.Fair}); this is exactly $h\in S$
(\texttt{fair\_iff\_mem\_feasibilitySet}).
\begin{theorem}[\texttt{feasibilitySet\_nonempty\_of\_fair}]
Fair information is consistent: if the truth satisfies every piece of
information then $S\ne\emptyset$.
\end{theorem}
Contrapositively, inconsistency is a certificate that some information is
wrong:
\begin{theorem}[\texttt{exists\_unfair\_of\_feasibilitySet\_eq\_empty}]
If $S=\emptyset$ then for every candidate $h$ there is an $i$ with
$h\notin S_i$.
\end{theorem}

\paragraph{Ideal information.}
The family is \emph{ideal} for $h$ when $S=\{h\}$ (\texttt{STE.Ideal}).
\begin{theorem}[\texttt{Ideal.eq\_of\_mem},
\texttt{ideal\_iff\_fair\_and\_subsingleton}]
Under ideal information every estimate equals the truth; and idealness is
equivalent to fairness together with feasibility being a subsingleton.
\end{theorem}

\paragraph{Topology (Carlson's setting).}
\begin{theorem}[\texttt{isClosed\_feasibilitySet},
\texttt{feasibilitySet\_nonempty\_of\_compact}]
If each $S_i$ is closed then $S$ is closed. If moreover $\Xi$ is compact
and every finite subfamily is consistent (the finite intersection
property, \texttt{FiniteConsistency}), then $S\ne\emptyset$.
\end{theorem}
\noindent The compactness result is the abstract reason STE is well posed
on Carlson's finite key spaces (discrete $\Rightarrow$ compact;
\texttt{isClosed\_feasibilitySet\_of\_discrete}) and on bounded closed
convex constraint sets.

\paragraph{Convexity (Combettes's Hilbert-space setting).}
\begin{theorem}[\texttt{convex\_feasibilitySet}]
If each $S_i$ is convex then $S$ (and every partial feasibility set) is
convex.
\end{theorem}
\noindent This is the geometric precondition for the projection
algorithms of \citet{youla1982,bauschke1996}, a planned extension.

\paragraph{STE decryption skeleton (Carlson).}
A \texttt{CipherSystem} is a keyed family of injective encryption maps.
Observing a ciphertext $c$ induces the key property set
$\{k \mid \exists m\in\mathrm{Meaning},\ \mathrm{enc}\,k\,m=c\}$; the
feasible key set is the STE feasibility set of these. The cipher module
now lives in the \texttt{STE.Carlson} namespace.
\begin{theorem}[\texttt{STE.Carlson.CipherSystem.fair\_of\_genuine},
\texttt{STE.Carlson.CipherSystem.feasibleKeys\_nonempty\_of\_genuine}]
Genuine traffic (ciphertexts that really are encryptions of meaningful
messages under a true key $k_0$) is fair for $k_0$; hence the feasible
key set is nonempty.
\end{theorem}
\begin{theorem}[\texttt{STE.Carlson.CipherSystem.Unicity.eq\_of\_feasible}]
At unicity (feasible key set $=\{k_0\}$, i.e.\ ideal information) every
feasible key is the true key: decryption succeeds by feasibility alone.
\end{theorem}

\paragraph{Cipher key counting (Carlson, Lemma~4.1; new).}
Model a substitution cipher over a finite alphabet $A$ by its key space
$\mathrm{Perm}\,A$. Given a message whose distinct symbols form
$T\subseteq A$, two keys decrypt it identically exactly when they agree
on $T$; Carlson's Lemma~4.1 counts the keys equivalent to the true key.
\begin{theorem}[\texttt{STE.Carlson.card\_consistent\_keys}]
For a substitution cipher with true key $\sigma$ over $A$ and observed
symbol set $T$,
\[ \#\{\tau\in\mathrm{Perm}\,A : \forall a\in T,\ \tau a=\sigma a\}
   = (|A|-|T|)!. \]
\end{theorem}
\noindent \texttt{Ste/Carlson/Counting.lean} is deliberately independent
of the \texttt{CipherSystem} layer, so on its own this is a permutation
count, not yet a statement about \texttt{feasibleKeys}. The bridge is
made in \texttt{Ste/Carlson/SheafBridge.lean}, which imports both:
\texttt{STE.Carlson.mem\_keyPropertySet\_iff\_decConstraint} identifies
Carlson's STE property set on keys with the decryption constraint read
through the inverse key, and
\texttt{STE.Carlson.card\_restrict\_fiber} identifies the count above
with the cardinality of a fiber of the restriction map to the observed
symbols (the map defining
\texttt{Ste.VariablePresheaf.localSections}). Read through those two
lemmas, $(|A|-|T|)!$ is the residual STE key ambiguity after one
observation.
\noindent The proof reduces, by left-translating the constraint by
$\sigma^{-1}$ (an explicit bijection), to the fixing-set count
\texttt{STE.Carlson.card\_perm\_fixing}: the permutations fixing $T$
pointwise number $(|A|-|T|)!$, via
\texttt{Equiv.Perm.subtypeEquivSubtypePerm} and
\texttt{Fintype.card\_perm}. The total key count is
\texttt{STE.Carlson.card\_substitution\_keys}: $\#\mathrm{Perm}\,A=|A|!$.

\paragraph{Cipher reduction (Carlson, Theorem~5.2 / Lemma~5.1 /
Corollary~5.2 / Corollary~5.3; new).}
Model an $n$-bit permutation cipher by $\mathrm{Perm}(\mathrm{Fin}\,n)$
acting on the value space $\mathrm{Fin}\,n\to\mathrm{Bool}$: a position
permutation $\rho$ induces the value permutation
$\mathrm{coperm}\,\rho$, acting by
$(\mathrm{coperm}\,\rho)\,f\,x=f(\rho^{-1}x)$.
\begin{theorem}[\texttt{STE.Carlson.coperm\_apply},
\texttt{STE.Carlson.coperm\_mul},
\texttt{STE.Carlson.coperm\_injective},
\texttt{STE.Carlson.permutation\_reduces\_to\_substitution}]
$\mathrm{coperm}$ acts by relabelling coordinates,
$(\mathrm{coperm}\,\rho)\,f\,x=f(\rho^{-1}x)$; it is a group
homomorphism, $\mathrm{coperm}(\rho_1\rho_2)=
\mathrm{coperm}\,\rho_1\cdot\mathrm{coperm}\,\rho_2$ (packaged as
\texttt{copermHom}); and it is \emph{injective}. Hence the
permutation-cipher keys form a subgroup of the substitution-cipher keys
(Theorem~5.2). The Lean statement of
\texttt{permutation\_reduces\_to\_substitution} is the injectivity of
\texttt{copermHom}: injectivity of an explicit homomorphism, not the
existence of an image, is the reduction's content.
\end{theorem}
\begin{theorem}[\texttt{STE.Carlson.card\_permutation\_lt\_card\_substitution},
\texttt{STE.Carlson.substitution\_not\_reducible\_to\_permutation}]
For an $n$-bit block with $n\ge 2$ there are strictly more substitution
keys than permutation keys ($n! < (2^n)!$), so \emph{no} injective
realization of permutation keys as substitution keys is onto: some
substitution cipher is not a permutation cipher (Lemma~5.1 /
Corollary~5.2). The non-surjectivity is quantified over every faithful
embedding, so it does not depend on a particular construction.
\end{theorem}
\begin{theorem}[\texttt{STE.Carlson.psp\_apply},
\texttt{STE.Carlson.psp\_reduces\_to\_substitution}]
For \emph{independent} position permutations $\rho_1,\rho_2$ and a
substitution key $s$, the composite
$\mathrm{coperm}\,\rho_1\cdot s\cdot\mathrm{coperm}\,\rho_2$ acts by
$f\mapsto \bigl(x\mapsto s(f\circ\rho_2^{-1})(\rho_1^{-1}x)\bigr)$, and
the set $\mathrm{pspKeys}\,n$ of keys so realized is \emph{exactly}
$\mathrm{Perm}(\mathrm{Fin}\,n\to\mathrm{Bool})$. The inclusion
$\subseteq$ is closure (a PSP cipher is a single substitution cipher);
the reverse inclusion, taking $\rho_1=\rho_2=1$, says the PSP
construction gains no keys, so the reduction is an equality of key
spaces (Corollary~5.3).
\end{theorem}
\noindent Throughout, the numbers 5.1/5.2/5.3 are used as attribution
labels for Carlson's Chapter~5 reduction chain (Theorem~5.2 for
$P\Rightarrow S$; Lemma~5.1 / Corollary~5.2 for the failure of
$S\Rightarrow P$; Corollary~5.3 for $PSP\Rightarrow S$), applied
consistently here and in the Lean sources. The copy of the dissertation
we hold is a scanned PDF whose body text is not machine-readable, so the
numbering could not be re-verified mechanically; each Lean statement
stands on its own proof regardless.

\paragraph{Asymptotic unicity (\texttt{Ste/Carlson/Asymptotic.lean}).}
Carlson's asymptotic claim is that, as intercepted ciphertexts
accumulate, the feasible key set collapses onto the true key. The file
mechanizes the deterministic core. Write
$\mathrm{partialKeys}\,n$ for the partial STE feasibility set of the key
property sets induced by the first $n$ ciphertexts of a stream
$\mathrm{obs}:\mathbb N\to C$, and define
$\mathrm{AsymptoticUnicity}$ to hold for $k_0$ when
$\mathrm{partialKeys}\,n=\{k_0\}$ eventually along
$\mathrm{atTop}$.
\begin{theorem}[\texttt{STE.Carlson.CipherSystem.partialKeys\_antitone},
\texttt{...asymptoticUnicity\_iff\_tendsto},
\texttt{...iInter\_partialKeys}]
$n\mapsto\mathrm{partialKeys}\,n$ is antitone (more traffic can only
narrow the key search); asymptotic unicity is exactly convergence of
$\mathrm{partialKeys}$ to $\{k_0\}$ in the discrete topology on
$\mathcal P(K)$; and $\bigcap_n \mathrm{partialKeys}\,n$ is the feasible
key set of the whole stream.
\end{theorem}
\begin{theorem}[\texttt{...asymptoticUnicity\_iff},
\texttt{...asymptoticUnicity\_iff\_unicity}]
For a finite key space, asymptotic unicity for $k_0$ is equivalent to
$\bigcap_n\mathrm{partialKeys}\,n=\{k_0\}$, hence to
$\mathrm{Unicity}$ for the entire observation stream. The proof rests on
\texttt{STE.antitone\_eventually\_eq\_iInter}: over a finite space an
antitone sequence of sets is eventually equal to its intersection.
\end{theorem}
\noindent Scope note, stated honestly: Carlson's dissertation asserts the
asymptotic result \emph{almost surely} with respect to a source
distribution on plaintexts. Nothing measure-theoretic is attempted here;
everything above is deterministic, and the a.s.\ version remains future
work (it needs a probability space on message streams from
\texttt{Mathlib.MeasureTheory}).

\paragraph{Module inventory.}
The Carlson development is \texttt{Ste/Carlson/}:
\texttt{Cipher.lean} (STE decryption skeleton),
\texttt{Counting.lean} (Lemma~4.1),
\texttt{Reduction.lean} (the Chapter~5 chain),
\texttt{Asymptotic.lean} (above),
\texttt{SheafBridge.lean} (the bridge to the support / variable-presheaf
layer), and \texttt{SetOfSets.lean} (second-order sources contributing a
set of candidate possibility sets). All six are imported by the umbrella
module \texttt{Ste/Carlson.lean}. Adjacent to it,
\texttt{Ste/PluralFeasibility.lean} treats J\o sang-style hyper sources
in the crisp STE setting: a source offering a family $H_i$ of
alternative property sets carries the information $\bigcup_{S\in H_i}S$,
and the resulting plural feasibility set decomposes as the union, over
all selections of one branch per source, of the ordinary feasibility
sets.

\section{Verification method}
The environment used to author this work cannot download a Lean toolchain
(egress policy). We therefore treat the GitHub Action as the trusted
verifier: every push to \texttt{main} builds the \texttt{Ste} library
with \texttt{leanprover/lean-action} against the pinned Mathlib cache. A
result is considered established only once its build is green. This
operationalizes the project's rule to distrust results that are not
machine proven.

\section{Future work}
With Lemma~4.1, the reduction chain (Theorem~5.2, Lemma~5.1 /
Corollary~5.2, Corollary~5.3) and the deterministic asymptotic-unicity
core now discharged, the remaining targets, in increasing
depth, are: the almost-sure form of asymptotic unicity, requiring a
probability space on plaintext streams
(\texttt{Mathlib.MeasureTheory}) --- the one gap explicitly flagged in
\texttt{Ste/Carlson/Asymptotic.lean}; Carlson's general counting
Theorem~4.1 (the
$\binom{|S_t|}{C}(|B|-|S_t|)!$ count with a chosen-symbol constraint);
POCS convergence over Hilbert spaces; and Theorem~2.3, that property
sets are AEP typical sets, placing information theory under STE. Targets
and their open/proven status are tracked in the project \texttt{log/}.

\bibliographystyle{plainnat}
\bibliography{../../refs}

\end{document}
